\section{Results}
\label{sec:results}

\subsection{The global execution horizon matters}

On task 0, global success is 29/50, 31/50, 42/50, 45/50, and 42/50 for
$h=1,2,4,8,16$, respectively.  Thus neither the most frequent replanning nor
the longest evaluated open-loop execution is best.  Policy-query rate decreases
from 1.000 at $h=1$ to 0.065 at $h=16$, making success and policy-inference
demand separate
deployment axes.

\subsection{Horizon assignment to action groups matters}

Fig.~\ref{fig:results}(a) shows the complete task-0 group landscape.  The best
evaluated point is $(4,16)$ with 47/50 success.  More importantly than this
single maximum, performance is asymmetric about the diagonal.  For every one of
the ten distinct horizon pairs, assigning the longer value to the gripper
outperforms the reversed assignment.  Differences range from 0.08 for $(1,2)$
versus $(2,1)$ to 0.28 for both $(1,8)$ versus $(8,1)$ and $(2,16)$ versus
$(16,2)$.  This does not establish a universal longer-gripper rule.  It shows
that the same horizon values and essentially the same query cadence can yield
different outcomes when assigned to different action groups.

\begin{figure*}[t]
    \centering
    \includegraphics[width=0.93\textwidth]{results_overview.pdf}
    \caption{Static horizon structure.  (a) Complete task-0 success landscape;
    diagonal cells are global-equivalent controls.  (b) Best evaluated global
    and off-diagonal success for each of tasks 1--9.  Panel (b) is a retrospective
    per-task oracle-style diagnostic, not the performance of one deployable
    selector.}
    \label{fig:results}
\end{figure*}

\subsection{Improvements persist at controlled query rates}

Table~\ref{tab:budget} isolates three pairs whose policy-query rates are nearly
identical.  On task 0, global $h=4$ and $(4,16)$ both query at approximately
0.252 per step.  The group-specific configuration gains five paired successes
with no reverse discordance; the exact two-sided diagnostic is $p=0.0625$, which
we do not describe as statistically significant.  Tasks 1 and 4 show larger
paired differences at approximately 0.50 queries per step.  These examples rule
out increased policy-query frequency as the explanation for those particular
gains, but do not identify their causal mechanism.

\begin{table}[t]
\caption{Budget-controlled paired comparisons}
\label{tab:budget}
\centering
\scriptsize
\setlength{\tabcolsep}{2.4pt}
\begin{tabular}{clccc}
\toprule
Task & Comparison & Success & Query rate & Discord. / $p$ \\
\midrule
0 & $h=4$ vs. $(4,16)$ & 42/50 $\to$ 47/50 & .252/.252 & 5--0 / .0625 \\
1 & $h=2$ vs. $(2,16)$ & 6/20 $\to$ 14/20 & .500/.501 & 8--0 / .0078 \\
4 & $h=2$ vs. $(2,16)$ & 10/20 $\to$ 17/20 & .501/.502 & 7--0 / .0156 \\
\bottomrule
\end{tabular}
\vspace{-1mm}
\begin{flushleft}\footnotesize
`Discord.' gives group-specific-only versus global-only successes.  $p$ is an
uncorrected exact two-sided McNemar diagnostic.
\end{flushleft}
\end{table}

\subsection{Off-diagonal structure recurs across tasks}

Across tasks 1--9, the retrospective per-task best global success averages
0.700, while the best evaluated off-diagonal point averages 0.761
[Fig.~\ref{fig:results}(b), Table~\ref{tab:cross_task}].  This $+0.061$ is a static oracle-style diagnostic:
it selects a configuration after observing each task's results and is not a
deployable method.  The best evaluated class is strictly off-diagonal on six
tasks and tied between diagonal and off-diagonal on three; none is diagonal-only.
At least one off-diagonal point strictly improves a global point at no higher
measured query rate on 6/9 tasks.

Directionality is recurring but not universal.  Across 36 symmetric comparisons
on these tasks, the longer-gripper assignment wins 24, the reversed assignment
wins 6, and 6 tie.  The reversals are important: group identity matters, but a
single directional explanation is not supported.

\begin{table*}[t]
\caption{Cross-task retrospective best evaluated static configurations}
\label{tab:cross_task}
\centering
\scriptsize
\setlength{\tabcolsep}{5.0pt}
\begin{tabular}{cclclc}
\toprule
Task & Best global horizon(s) & Global success & Best off-diagonal pair(s) & Off-diagonal success & Difference \\
\midrule
1 & 4, 16 & .50 & $(2,16)$ & .70 & +.20 \\
2 & 4 & .75 & $(2,16)$, $(4,16)$ & .80 & +.05 \\
3 & 16 & .90 & $(16,8)$ & .90 & .00 \\
4 & 4, 16 & .80 & $(2,16)$ & .85 & +.05 \\
5 & 16 & .85 & $(16,8)$ & .90 & +.05 \\
6 & 16 & .65 & $(8,16)$ & .65 & .00 \\
7 & 8, 16 & .55 & $(4,16)$ & .65 & +.10 \\
8 & 2, 4 & .40 & $(2,8)$, $(8,2)$ & .40 & .00 \\
9 & 2, 4, 8 & .90 & $(16,8)$, $(4,16)$ & 1.00 & +.10 \\
\midrule
Macro & --- & .700 & per-task selection & .761 & +.061 \\
\bottomrule
\end{tabular}
\end{table*}

\subsection{Preferred static settings vary, but a common pair is strong}

The best evaluated pairs differ by task in the coarse grid.  A separate post-hoc
common-set check across all ten tasks nevertheless selects $(4,16)$ as the best
single group-specific pair, with macro success 0.734; the best global point is
$h=16$ at 0.699.  Their query rates differ (0.252 versus 0.065), so this is not a
budget-matched superiority claim.  The retrospective per-task group oracle over
the same common set reaches 0.779, a 0.045 gap, while leave-one-task-out static
selection chooses $(4,16)$ in every split.  Current evidence therefore supports
a strong universal static baseline and leaves the need for dynamic scheduling
open.

\subsection{Exploratory execution-trace context}

The step logs permit a descriptive, non-causal check of the executed controls.
Across all 50 task-0 rollouts, mean absolute gripper step change is 0.090 for
global $h=4$, 0.053 for global $h=8$, 0.048 for $(4,16)$, and 0.067 for
$(16,4)$.  Mean arm $\ell_2$ step change is 0.085, 0.077, 0.082, and 0.085 in
the same order; mean gripper sign changes per episode are 8.12, 3.84, 3.64, and
5.74.  These aggregates mix successful and failed trajectories of different
lengths.  They show that the executor changes realized action traces, but they
do not show that reduced gripper variation causes success or isolate one
mechanism for the horizon preference.
